% =============================================================================
%  lesson-presentation.tex — عرض درس تقنية المعلومات (الدرس 01)
%  المجال 1: بيئة التعامل مع الحاسوب — الوحدة 1: تقنية المعلومات (IT)
%  السنة الدراسية: 2026 / 2027
% =============================================================================
%  البناء (من هذا المجلد):
%    xelatex -interaction=nonstopmode -halt-on-error -output-directory=build lesson-presentation.tex
% =============================================================================

\documentclass[aspectratio=169, 11pt]{beamer}

% ── النواة المشتركة للعروض (الحزم + الهوية + المظهر + الخلاصات) ─────────────
\input{"../../shared/presentation.tex"}

% ══════════════════════════════════════════════════════════════════════════════
\title{تقنية المعلومات}
\subtitle{المجال 1: بيئة التعامل مع الحاسوب — الوحدة 1: تقنية المعلومات (IT)}
\author{الأستاذ لفقير عبدالجليل}
\institute{ثانوية الشهيد حكومي العيد أدرار}
\date{السنة الدراسية 2027/2026}

% ══════════════════════════════════════════════════════════════════════════════
\begin{document}

% ── شريحة العنوان: الترويسة الرسمية ─────────────────────────────────────────
\begin{frame}[plain]
  \officialheader
  \vfill
  \begin{center}
    {\color{primary}\Huge\bfseries\inserttitle}
  \end{center}
  \vfill
  \begin{center}
    {\color{rulegray}\footnotesize \insertauthor\ — \insertdate}
  \end{center}
\end{frame}

% ── أهداف الدرس ──────────────────────────────────────────────────────────────
\begin{frame}{أهداف الدرس}
  \begin{itemize}
    \item يتعرّف التلميذ على مفاهيم \hl{التقنية}، \hl{التكنولوجيا}، \hl{الإعلام}، و\hl{الاتصال}.
    \item يتعرّف على مفهوم \hl{تكنولوجيا الإعلام والاتصال (TIC)} ومكوّناتها.
    \item يميّز بين \hl{البيانات} و\hl{المعلومات} وعملية \hl{المعالجة} الآلية.
    \item يتعرّف على مفهوم \hl{المعلوماتية (Informatique)} وأصل الكلمة.
    \item يتعرّف على \hl{مراحل تطور تقنية المعلومات}: الأجهزة والبرامج.
  \end{itemize}
\end{frame}

% ── المدخل: المصطلحات ومجالات استعمالها ─────────────────────────────────────
\begin{frame}{المصطلحات ومجالات استعمالها}
  \centering
  {\footnotesize
  \begin{tabularx}{\textwidth}{@{}|>{\raggedright\arraybackslash}p{3.4cm}|>{\raggedright\arraybackslash}X|@{}}
    \hline
    \rowcolor{primary}
    \colorbox{primary}{\parbox{3.2cm}{\centering\color{white}\textbf{المصطلح}}} &
    \cellcolor{primary}\color{white}\textbf{مجالات الاستعمال} \\
    \hline
    \textbf{\color{primary}المعلوماتية / الإعلام الآلي} &
    إنجاز الدروس والواجبات، متابعة السلع والفواتير في المحلات، ومعالجة الطلبات والمسابقات وتسجيل التلاميذ لشهادة التعليم المتوسط (وزارة التربية). \\
    \hline
    \textbf{\color{primary}التكنولوجيا} &
    البيوت الذكية (مثل Siri)، الروبوتات، الزراعة الآلية، والطبّ عن بُعد. \\
    \hline
    \textbf{\color{primary}الإعلام والاتصال} &
    الإذاعة والبثّ (عتاد وميكروفون وكاميرا وتطبيقات المونتاج)، والبثوث المباشرة على منصات التواصل. \\
    \hline
    \rowcolor{lightbg}
    \textbf{القاسم المشترك} & جميعها \key{تُيسِّر حياة الإنسان}. \\
    \hline
  \end{tabularx}\par}
\end{frame}

% ── الإشكالية ────────────────────────────────────────────────────────────────
\begin{frame}{الإشكالية}
  \begin{block}{وضعية انطلاق}
    نلاحظ أنّ «المعلوماتية» و«الإعلام الآلي» يدلّان على \key{المعنى نفسه}،
    أمّا «التكنولوجيا» و«الإعلام والاتصال» فـ\hl{مفاهيم مختلفة}.\\[4pt]
    \textbf{هل تعبّر هذه الكلمات عن مفهوم واحد؟}\\[4pt]
    فما \textbf{مفهوم تقنية المعلومات}؟ وما \textbf{مراحل تطورها}؟
  \end{block}
\end{frame}

% ── 1. المفاهيم الأولية ─────────────────────────────────────────────────────
\sectionframe{1}{المفاهيم الأولية}

\begin{frame}{التقنية والتكنولوجيا}
  \begin{block}{التقنية (Technique)}
    هي \key{المهارة والقدرة الفردية} على ممارسة الأعمال والأنشطة لإنجاز منتج معيّن.
  \end{block}
  \begin{block}{التكنولوجيا (Technologie)}
    كلمة ذات \key{أصل يوناني} مركّبة من مقطعين:
    \centerline{\key{Techno} (مهارة فنية) \;+\; \key{Logos} (علم)}
    وتعني \hl{علم المهارات والتقنيات التطبيقية} والاستغلال المنطقي للمعارف العلمية.
  \end{block}
\end{frame}

\begin{frame}{الإعلام والاتصال}
  \begin{block}{الإعلام (Information)}
    هو تقديم الأخبار والمعلومات و\key{نقل رسالة} من \hl{مرسل} إلى \hl{مستقبل}.
  \end{block}
  \begin{block}{الاتصال (Communication)}
    هو عملية \key{التفاعل ونقل الأفكار والمعلومات} بين مرسل ومستقبل في سياق معيّن.
  \end{block}
  \vspace{4pt}
  \begin{center}
  \begin{tikzpicture}[
    node distance=2.6cm,
    box/.style={draw=primary, fill=lightbg, thick, rounded corners=3pt,
                minimum width=2.6cm, minimum height=0.9cm, align=center},
    ar/.style={-{Latex[length=3mm]}, thick, primary}
  ]
    \node[box] (s) {مرسل};
    \node[box, right=of s] (m) {رسالة};
    \node[box, right=of m] (r) {مستقبل};
    \draw[ar] (s) -- (m);
    \draw[ar] (m) -- (r);
  \end{tikzpicture}
  \end{center}
\end{frame}

% ── 2. مفهوم TIC ────────────────────────────────────────────────────────────
\sectionframe{2}{تكنولوجيا الإعلام والاتصال (TIC)}

\begin{frame}{مفهوم TIC}
  \begin{block}{التعريف}
    هي \key{التلاقي والتزاوج} بين عتاد الحواسيب، البرمجيات، وشبكات الاتصال
    لـ\hl{جمع البيانات} وتخزينها و\hl{معالجتها ونقلها}.
  \end{block}
  \vspace{6pt}
  \begin{center}
  \begin{tabular}{@{}c@{\hspace{0.6cm}}c@{\hspace{0.6cm}}c@{}}
    \colorbox{lightbg}{\parbox{3cm}{\centering\key{عتاد الحواسيب}}} &
    \colorbox{lightbg}{\parbox{3cm}{\centering\key{البرمجيات}}} &
    \colorbox{lightbg}{\parbox{3cm}{\centering\key{شبكات الاتصال}}} \\
  \end{tabular}
  \end{center}
\end{frame}

\begin{frame}{مكوّنات TIC الأساسية}
  \begin{block}{1. تقنيات المعالجة والتخزين}
    \begin{itemize}
      \item \key{الأجهزة:} كمبيوتر، هواتف.
      \item \key{البرامج:} أنظمة تشغيل، تطبيقات.
    \end{itemize}
  \end{block}
  \vspace{6pt}
  \begin{block}{2. شبكات الاتصال}
    \begin{itemize}
      \item السلكية.
      \item اللاسلكية.
      \item السمعي البصري.
    \end{itemize}
  \end{block}
\end{frame}

% ── 3. البيانات والمعلومات ──────────────────────────────────────────────────
\sectionframe{3}{البيانات والمعلومات والمعالجة}

\begin{frame}{البيانات والمعلومات والمعالجة}
  \begin{description}
    \item[\key{البيانات (Data):}] الحقائق والأرقام والرموز \hl{الأولية الخام}.
    \item[\key{المعلومات (Informations):}] البيانات التي جرت \hl{معالجتها وتنظيمها} فأصبحت ذات معنى.
    \item[\key{المعالجة (Processing):}] العمليات \hl{الحسابية والمنطقية} لتحويل البيانات إلى معلومات.
  \end{description}
\end{frame}

\begin{frame}{دورة معالجة البيانات}
  \begin{center}
  \begin{tikzpicture}[
    node distance=2.3cm,
    box/.style={draw=primary, fill=lightbg, thick, rounded corners=4pt,
                minimum width=3.0cm, minimum height=1.3cm, align=center},
    ar/.style={-{Latex[length=4mm]}, line width=1.1pt, primary}
  ]
    \node[box] (d) {\key{بيانات أولية}\\[2pt]\footnotesize (Data)};
    \node[box, right=of d] (p) {\key{معالجة آلية}\\[2pt]\footnotesize (Processing)};
    \node[box, right=of p] (i) {\key{معلومات مفيدة}\\[2pt]\footnotesize (Information)};
    \draw[ar] (d) -- (p);
    \draw[ar] (p) -- (i);
  \end{tikzpicture}
  \end{center}
  \vspace{6pt}
  \begin{block}{مثال}
    \hl{[علي — 2008 — 14.5]} \quad$\longrightarrow$\quad
    \hl{ترتيب وفرز} \quad$\longrightarrow$\quad
    \key{علي: ناجح بمعدل 14.5}
  \end{block}
\end{frame}

% ── 4. المعلوماتية ──────────────────────────────────────────────────────────
\sectionframe{4}{مفهوم المعلوماتية (Informatique)}

\begin{frame}{أصل كلمة المعلوماتية}
  \begin{block}{التركيب اللغوي}
    كلمة فرنسية مركّبة من جزأين:
  \end{block}
  \vspace{4pt}
  \begin{center}
  \begin{tikzpicture}[
    node distance=1.2cm,
    part/.style={draw=secondary, fill=lightbg, thick, rounded corners=3pt,
                 minimum width=4.2cm, minimum height=1.0cm, align=center},
    res/.style={draw=primary, fill=lightbg, thick, rounded corners=3pt,
                minimum width=8.6cm, minimum height=1.1cm, align=center}
  ]
    \node[part] (a) {\key{INFORMATION}\\[1pt]\footnotesize (معلومات)};
    \node[part, right=of a] (b) {\key{AUTOMATIQUE}\\[1pt]\footnotesize (آلية)};
    \node[res, below=1.0cm of a, xshift=2.7cm] (r)
      {\Large\key{INFORMATIQUE} \;\footnotesize (المعلوماتية)};
    \draw[-{Latex[length=3mm]}, thick, primary] (a.south) -- (r.north west);
    \draw[-{Latex[length=3mm]}, thick, primary] (b.south) -- (r.north east);
  \end{tikzpicture}
  \end{center}
\end{frame}

\begin{frame}{تعريف المعلوماتية}
  \begin{block}{التعريف}
    هي \key{علم المعالجة الآلية للمعلومات} باستخدام جهاز الحاسوب والبرمجيات المخصّصة لذلك.
  \end{block}
\end{frame}

% ── 5. مراحل تطور تقنية المعلومات ────────────────────────────────────────────
\sectionframe{5}{مراحل تطور تقنية المعلومات}

\begin{frame}{تطور أجهزة المعالجة}
  \begin{center}
  \begin{tikzpicture}[
    node distance=1.0cm,
    gen/.style={draw=primary, fill=lightbg, thick, rounded corners=3pt,
                minimum width=3.1cm, minimum height=1.1cm, align=center},
    ar/.style={-{Latex[length=3mm]}, line width=1.1pt, primary}
  ]
    \node[gen] (a) {حاسوب بحجم \key{غرفة}};
    \node[gen, left=of a] (b) {حاسوب \key{محمول}};
    \node[gen, left=of b] (c) {جهاز \key{لوحي}};
    \node[gen, left=of c] (d) {هاتف \key{ذكي}};
    \draw[ar] (a) -- (b);
    \draw[ar] (b) -- (c);
    \draw[ar] (c) -- (d);
  \end{tikzpicture}
  \end{center}
  \vspace{4pt}
  \begin{block}{ومع ذلك}
    يشمل التطوّر أيضًا \key{مكوّنات الحاسوب}: الرامات وبطاقة الشاشة.
  \end{block}
\end{frame}

\begin{frame}{تطور البرامج: أنظمة التشغيل}
  \centering
  \begin{tabular}{@{}c@{\hspace{0.7cm}}c@{\hspace{0.7cm}}c@{}}
    \colorbox{lightbg}{\parbox{3.4cm}{\centering\key{Windows}}} &
    \colorbox{lightbg}{\parbox{3.4cm}{\centering\key{Linux}}} &
    \colorbox{lightbg}{\parbox{3.4cm}{\centering\key{macOS}}} \\
  \end{tabular}
  \vspace{6pt}
  \begin{block}{مسار تطوّر Windows}
    {\footnotesize MS-DOS ثمّ 1.0 و2.0 و3.0، ثمّ 95 و98، فـ XP وVista، ثمّ 7 و8، ثمّ 10 و11.}
  \end{block}
  \vspace{4pt}
  \begin{block}{أشهر توزيعات Linux}
    {\footnotesize Ubuntu وArch Linux وRed Hat وFedora.}
  \end{block}
\end{frame}

\begin{frame}{تطور التطبيقات}
  \begin{block}{البرامج التطبيقية}
    \begin{itemize}
      \item تطبيقات \hl{Adobe}.
      \item المتصفحات والألعاب.
      \item برامج \hl{Word} و\hl{Excel} و\hl{PowerPoint}.
    \end{itemize}
  \end{block}
  \vspace{6pt}
  \begin{block}{قاعدة التطور}
    كلُّ تحديثٍ لهذه التطبيقات يُعدّ \key{مرحلةً} من مراحل تطور تقنية المعلومات.
  \end{block}
\end{frame}

% ── الخلاصة ──────────────────────────────────────────────────────────────────
\begin{frame}{الخلاصة}
  \begin{block}{ما يجب تذكّره}
    \begin{itemize}
      \item \key{التقنية} مهارة فردية، و\key{التكنولوجيا} علم تطبيقي.
      \item \key{TIC} = تلاقي العتاد + البرمجيات + شبكات الاتصال.
      \item \hl{بيانات} $\xrightarrow{\ \text{معالجة}\ }$ \hl{معلومات}.
      \item \key{INFORMATIQUE} = INFORMATION + AUTOMATIQUE.
      \item تطوّر تقنية المعلومات = \hl{أجهزة المعالجة} + \hl{البرامج} (أنظمة التشغيل والتطبيقات).
    \end{itemize}
  \end{block}
\end{frame}

\begin{frame}[plain]
  \vfill
  \begin{center}
    {\color{primary}\Huge\bfseries شكرًا لحسن انتباهكم}
  \end{center}
  \vfill
\end{frame}

\end{document}
